\exercisesfor{1}

\begin{enumerate}
\def\labelenumi{\arabic{enumi}.}
\item
  设 \(\mathfrak{g}\) 是李代数，\(\mathfrak{a}\) 和 \(\mathfrak{b}\) 是 \(\mathfrak{g}\) 的理想（分别为特征理想）。则，集合中的 \(x \in \mathfrak{g}\) 若满足 \([x, \mathfrak{b}] \subset \mathfrak{a}\)，就是 \(\mathfrak{g}\) 的理想（分别为特征理想），称为把 \(\mathfrak{b}\) 传送到 \(\mathfrak{a}\) 中的传送子。证明 \(\mathcal{C}_{p+1}\mathfrak{g}\) 是把 \(\mathfrak{g}\) 传送到 \(\mathcal{C}_p\mathfrak{g}\) 中的传送子。
\item
  若 \(\mathfrak{g}\) 是 \(n\) 维、定义在 \(\mathbf{K}\) 上的李代数，且 \(\mathfrak{g}\) 的中心维数 \(\geqslant n - 1\)，则 \(\mathfrak{g}\) 是交换的。
\item
  设 M 是域 K 上不必结合的代数，\((e_{i})_{i\in I}\) 是向量空间 M 的一组基，\(c_{ijk}\) 是关于基 \((e_{i})\) 的 M 的结构常数。M 成为李代数的充要条件是 \(c_{ijk}\) 满足以下条件：（a）\(c_{iik}=0\)；（b）\(c_{ijk}=-c_{jik}\)；（c）\(\sum_{r\in I}(c_{ijr}c_{rkl}+c_{jkr}c_{ril}+c_{kir}c_{rjl})=0\) 对 I 中所有 i,j,k,l 成立。
\item
  \begin{enumerate}
  \def\labelenumii{(\alph{enumii})}
  \tightlist
  \item
    假设 2 在 K 中可逆。设 \(\mathfrak{a}\) 是 K 上的李代数。在 K-模 \(\mathfrak{a}' = \mathfrak{a} \times \mathfrak{a}\) 上按下式定义括号运算：
  \end{enumerate}
\end{enumerate}

\[
[ (x _ {1}, x _ {2}), (y _ {1}, y _ {2}) ] = ([ x _ {1}, y _ {1} ] + [ x _ {2}, y _ {2} ], [ x _ {1}, y _ {2} ] + [ x _ {2}, y _ {1} ]).
\]

    于是 \(\mathfrak{a}'\) 是李 K-代数，而映射 \(x\mapsto \frac{1}{2} (x,x)\)、\(x\mapsto \frac{1}{2} (x, - x)\) 是把 \(\mathfrak{a}\) 同构到理想 \(\mathfrak{b}\) 和 \(\mathfrak{c}\) 的同构；这些理想是 \(\mathfrak{a}'\) 的理想，且 \(\mathfrak{a}'\) 是它们的直和。（考虑 \(\mathrm{K}'\)，它是 \(\mathbf{K}\) 的二次扩域，以 1 和 \(k\) 为基，且 \(k^2 = 1\)。构造 \(\mathfrak{a}_{(\mathrm{K}',1)}\)，然后把标量环限制为 \(\mathbf{K}\)。注意，\(\frac{1}{2} (1 + k)\) 和 \(\frac{1}{2}(1 - k)\) 是 \(\mathrm{K}'\) 的幂等元，并且 \((1 + k)(1 - k) = 0\)。）

\begin{enumerate}
\def\labelenumi{(\alph{enumi})}
\setcounter{enumi}{1}
\tightlist
\item
    设 \(a\) 是实李代数，\(\mathfrak{g} = a_{(C)}\)，而 \(b\) 是把 \(\mathfrak{g}\) 的标量域限制为 \(\mathbf{R}\) 后得到的实李代数。证明 \(b_{(C)}\) 同构于 \(a_{(C)}'\)，其中 \(a'\) 是（a）中引入的代数。（考虑 \(\mathrm{V} = \mathbf{C} \otimes_{\mathbf{R}} \mathbf{C}\)，把它作为 \(\mathbf{C}\) 上的向量空间，并规定标量乘法为 \((z \otimes z')z'' = z \otimes z'z''\)。注意，\(\mathrm{V}\) 是实向量子空间 \(\mathrm{V}\) 的复化，其中该子空间由 \(1 \otimes 1\) 和 \(i \otimes i\) 生成；该子空间可识别为 \(\mathbf{R}\) 的二次扩域，其基为 1 和 \(k\)，且 \(k^2 = 1\)。）
\end{enumerate}

\begin{enumerate}
\def\labelenumi{\arabic{enumi}.}
\setcounter{enumi}{4}
\item
  设 \(\mathrm{V}\)、\(\mathrm{W}\) 是域 \(n + 1\) 上维数分别为 \(n\) 和 \(\mathbf{K}\) 的向量空间。证明 \(\mathfrak{gl}(\mathbf{W})\) 同构于 \(\mathfrak{sl}(\mathbf{V})\) 的一个子代数。（可设 W 是 \(\mathrm{V}\) 的超平面。取 \(e \in \mathbb{V}\)，且 \(e \notin \mathbb{W}\)。对于 \(x \in \mathfrak{gl}(\mathbf{W})\)，令 \(y(x)\) 为 \(\mathfrak{gl}(\mathbf{V})\) 中延拓 \(x\) 且满足 \(y(e) = -\mathrm{Tr}(x).e\) 的元素。证明映射 \(x \mapsto y(x)\) 是把 \(\mathfrak{gl}(\mathbf{W})\) 同构到 \(\mathfrak{sl}(\mathbf{V})\) 的一个子代数上的同构。）
\item
  李代数 g 为结合代数的充要条件是 \(\mathcal{D}\mathfrak{g}\) 包含于 \(\mathfrak{g}\) 的中心。
\end{enumerate}

¶ 7. 设 A 是一个环（可以没有单位元），满足在 A 中 \(2a = 0\) 蕴含 \(a = 0\)。令 U 是 A 的子环，也是由 A 关联的李 Z-代数的理想。

\begin{enumerate}
\def\labelenumi{(\alph{enumi})}
\tightlist
\item
  若 \(x, y\) 属于 U，则 \((xy - yx)\mathrm{A} \subset \mathrm{U}\)（写成：
\end{enumerate}

\[
(x y - y x) s = x (y s) - (y s) x - y (x s - s x))
\]

并且 \(\mathrm{A}(xy - yx)\mathrm{A} \subset \mathrm{U}\)（写成：

\[
r (x y - y x) s = (x y - y x) s r + r [ (x y - y x) s ] - [ (x y - y x) s ] r).
\]

\begin{enumerate}
\def\labelenumi{(\alph{enumi})}
\setcounter{enumi}{1}
\item
  假设 U 是交换的。设 \(x \in \mathbf{U}\)，\(s \in \mathbf{A}\)，且 \(y = xs - sx\)。证明 \((yt)^3 = 0\) 对所有 \(t \in \mathbf{A}\) 都成立。（元素 \(x\) 与 \(xs - sx\) 及 \(xs^2 - s^2x\) 交换；由此推出 \(2(xs - sx)^2 = 0\)，从而 \(y^2 = 0\)。同理，\((yt - ty)^2 = 0\) 对所有 \(t \in \mathbf{A}\) 成立。）
\item
  此后假设 A 的双边理想只有 \(\{0\}\) 和 A，并且对 A 的每个非零元素 y，都存在 \(t \in A\) 使 yt 不是幂零的。证明：若 \(U \neq A\)，则 U 包含于 A 的中心。（利用 \(\langle a\rangle\) 证明 U 交换，再利用 \(\langle b\rangle\)。）
\item
  设 V 是与 A 关联的李 Z-代数的理想。则 V 包含于 A 的中心，或 \(V \supset [A, A]\)。（对 \(t \in A\)，使用满足 \([t, A] \subset V\) 的 t 的集合 U 应用（c）；使用恒等式 \([xy, z] = [x, yz] + [y, zx]\)。）
\end{enumerate}

\begin{enumerate}
\def\labelenumi{\arabic{enumi}.}
\setcounter{enumi}{7}
\tightlist
\item
  若 \(x_{1}, x_{2}, x_{3}, x_{4}\) 是李代数的 4 个元素，则
\end{enumerate}

\[
\begin{array}{r l} {[ [ [ x _ {1}, x _ {2} ], x _ {3} ], x _ {4} ] + [ [ [ x _ {2}, x _ {1} ], x _ {4} ], x _ {3} ]} & \\ & \quad + [ [ [ x _ {3}, x _ {4} ], x _ {1} ], x _ {2} ] + [ [ [ x _ {4}, x _ {3} ], x _ {2} ], x _ {1} ] = 0. \end{array}
\]

\begin{enumerate}
\def\labelenumi{\arabic{enumi}.}
\setcounter{enumi}{8}
\tightlist
\item
  设 \(\mathfrak{g}\) 是一个李代数，并且 \([[x,y],y] = 0\) 对所有 \(x\) 和 \(y\) 都成立。
\end{enumerate}

\begin{enumerate}
\def\labelenumi{(\alph{enumi})}
\item
  证明 \(3[[x,y],z] = 0\)；其中 \(x,y,z\) 中的变量取自 \(\mathfrak{g}\)。（注意，\([[x,y],z]\) 是 \((x,y,z)\) 的交错三线性函数，并应用 Jacobi 恒等式。）
\item
  类似地，利用（a）和练习 8，证明 \(\left[[[x,y],z],t\right]=0\)；其中变量 \(x,y,z,t\) 取自 \(\mathfrak{g}\)。
\end{enumerate}

\begin{enumerate}
\def\labelenumi{\arabic{enumi}.}
\setcounter{enumi}{9}
\tightlist
\item
  设 L 是结合代数，g 是与之关联的李代数。L 的每个导子都是 g 的导子，但反过来不成立。（考虑一个交换结合代数的恒等映射。）
\end{enumerate}
% semantic-fragments: legacy-input-seam
\relax{}
\begin{enumerate}
\def\labelenumi{\arabic{enumi}.}
\setcounter{enumi}{10}
\item
  设 \(\mathfrak{g}\) 为一个李代数，\(\mathfrak{h}\) 为 \(\mathfrak{g}\) 的理想，且 \(\mathcal{D}\mathfrak{h} = \mathfrak{h}\) 。证明 \(\mathfrak{h}\) 是 \(\mathfrak{g}\) 的特征理想。
\item
  设 g 为一个李代数。
\end{enumerate}

\begin{enumerate}
\def\labelenumi{(\alph{enumi})}
\item
  导子 D 与 g 的所有内导子交换的充要条件是 D 将 g 映入其中心。
\item
  以下设 \(\mathfrak{g}\) 的中心为零，于是将 \(\mathfrak{g}\) 视为由 \(\mathfrak{D}\) 的所有导子组成的李代数 \(\mathfrak{g}\) 的一个理想。证明 \(\mathfrak{g}\) 在 \(\mathfrak{D}\) 中的中心化子为零（使用 (a)）。特别地，\(\mathfrak{D}\) 的中心为零。
\item
  证明 \(\Delta\) 的满足 \(\mathfrak{D}\) 的导子 \(\Delta(\mathfrak{g}) = \{0\}\) 为零。 \(([\Delta(\mathfrak{D}),\mathfrak{g}]\subset \Delta([\mathfrak{D},\mathfrak{g}]) + [\mathfrak{D},\Delta(\mathfrak{g})]\subset \Delta(\mathfrak{g}) = \{0\}\) 因而由 (b) 得 \(\Delta(\mathfrak{D}) = \{0\}\)。)
\item
  进一步证明，若 \(\mathfrak{g} = \mathcal{D}\mathfrak{g}\)，则 \(\Delta\) 的每个导子 \(\mathfrak{D}\) 都是内导子（使用 (c) 和练习 11）。
\end{enumerate}

\begin{enumerate}
\def\labelenumi{\arabic{enumi}.}
\setcounter{enumi}{12}
\tightlist
\item
  设 \(\mathfrak{g}\) 为一个李代数，\(\mathfrak{a}\) 和 \(\mathfrak{b}\) 为 \(\mathfrak{g}\) 的两个子模。对于每个整数 \(i \geqslant 0\)，如下定义子模 \(\mathfrak{m}_i = \mathfrak{m}_i(\mathfrak{a},\mathfrak{b})\) 和 \(\mathfrak{n}_i = \mathfrak{n}_i(\mathfrak{a},\mathfrak{b})\)：\(\mathfrak{m}_1 = \mathfrak{b}, \mathfrak{m}_{i+1} = [\mathfrak{m}_i,\mathfrak{b}], \mathfrak{n}_1 = \mathfrak{a}, \mathfrak{n}_{i+1} = [\mathfrak{n}_i,\mathfrak{b}]\)。证明
\end{enumerate}

\[
[ a, m _ {i} ] \subset n _ {i + 1}
\]

对于 \(i = 1,2,\ldots\)（对 \(i\) 作归纳，注意到

\[
\mathfrak {m} _ {i} ([ a, b ], b) = \mathfrak {m} _ {i} (a, b)
\]

以及 \(n_{i}([a,b],b)=n_{i+1}(a,b)\)。

¶ 14. 设 \(\mathfrak{a}\) 为一个李代数，\(\mathfrak{b}\) 为 \(\mathfrak{a}\) 的子代数。连接 \(\mathfrak{a}\) 与 \(\mathfrak{b}\) 的合成列是一个递减序列 \((\mathfrak{a}_i)_{0\leqslant i\leqslant n}\)，其中各项是 \(\mathfrak{a}\) 的子代数，满足 \(\mathfrak{a}_0 = \mathfrak{a}\)、\(\mathfrak{a}_n = \mathfrak{b}\)，且 \(\mathfrak{a}_{i+1}\) 是 \(\mathfrak{a}_i\) 的理想，其中 \(0\leqslant i < n\)。称 \(\mathfrak{b}\) 为 \(\mathfrak{a}\) 的次不变子代数，当且仅当存在连接 \(\mathfrak{a}\) 与 \(\mathfrak{b}\) 的合成列。

\begin{enumerate}
\def\labelenumi{(\alph{enumi})}
\item
  设 \(\mathfrak{b}\) 是 \(\mathfrak{a}\) 的次不变子代数，\((\mathfrak{a}_i)_{0\leqslant i\leqslant n}\) 是连接 \(\mathfrak{a}\) 与 \(\mathfrak{b}\) 的合成列。由练习 13 推出 \([\mathcal{C}^{n+p}\mathfrak{b},\mathfrak{a}]\subset\mathcal{C}^{p}\mathfrak{b}\)，注意到当 \([\mathfrak{a}_i,\mathfrak{b}]\subset\mathfrak{a}_{i+1}\) 时 \(1\leqslant i<n\)。
\item
  由 (a) 推出，所有 \(\mathcal{C}^{\infty}\mathfrak{b}\)（\(\mathcal{C}^{p}\mathfrak{b}\)）的交 \(p = 0,1,2,\ldots\) 是 \(\mathfrak{a}\) 的理想。
\item
由 (b) 推出，若 \(\mathfrak{c}\) 是 \(\mathfrak{a}\) 的次不变子代数且满足 \(\mathfrak{c} = \mathcal{D}\mathfrak{c}\)，则它是 \(\mathfrak{a}\) 的理想。
\item
  当 K 是一个域且 \(\dim_{K}a<+\infty\) 时，证明所有 \(D^{\infty}b\) 的交 \(D^{p}b\left(p=0,1,2,\ldots\right)\) 是 a 的理想。（证明此交是 a 的次不变子代数，并应用 (c)。）
\end{enumerate}

¶ 15. (a) 设 a 为一个李代数，b 为 a 的子代数，c 为 b 的理想，z 为 c 在 a 中的中心化子。则 {[}b, b{]} ⊂ z。推出 b + z 是 a 的子代数，且 z 是其中的理想。

\begin{enumerate}
\def\labelenumi{(\alph{enumi})}
\setcounter{enumi}{1}
\item
  设 g 为一个李代数，a 为 g 的次不变子代数，b 为 a 的理想。若 b 在 a 中的中心化子以及 a 在 g 中的中心化子均为零，则 \(\mathfrak{z}\)（\(\mathfrak{b}\) 在 \(\mathfrak{g}\) 中的中心化子）为零。（令 \(\mathfrak{b} = \mathfrak{a} + \mathfrak{z}\)，由 \((a)\) 可知它是子代数；\(\mathfrak{a}\) 在 \(\mathfrak{b}\) 中次不变；若 \(\mathfrak{z} \neq \{0\}\)，则 \(\mathfrak{a} \neq \mathfrak{h}\)，因而 \(\mathfrak{a}\) 是 \(\mathfrak{a}'\) 的理想，其中 \(\mathfrak{a} \neq \mathfrak{a}' \subset \mathfrak{b}\)；取 \(a' \in \mathfrak{a}'\)、\(a' \notin \mathfrak{a}\)，并令 \(a' = a + z\)，其中 \(a \in \mathfrak{a}\)、\(z \in \mathfrak{z}\)、\(z \neq 0\)；证明 \([z, \mathfrak{a}]\) 包含于 \(\mathfrak{a}\) 且与 \(\mathfrak{b}\) 交换，于是 \([z, \mathfrak{a}] = \{0\}\)，矛盾。)
\item
设 \(\mathfrak{g}\) 为中心为零的李代数，\(\mathfrak{D}_1\) 为 \(\mathfrak{g}\) 的导子李代数，\(\mathfrak{D}_2\) 是 \(\mathfrak{D}_1, \ldots, \mathfrak{g}\) 的导子李代数，后者被识别为 \(\mathfrak{D}_1\) 的理想，\(\mathfrak{D}_1\) 被识别为 \(\mathfrak{D}_2, \ldots\) 的理想，……（参见练习 12 (b)）。利用 (b) 和练习 12 (b)，证明 \(\mathfrak{g}\) 在 \(\mathfrak{D}_i\) 中的中心化子对所有 \(i\) 均为零。（关于本练习的后续内容，参见 § 4，练习 15。）
\end{enumerate}

\begin{enumerate}
\def\labelenumi{\arabic{enumi}.}
\setcounter{enumi}{15}
\tightlist
\item
  设 g 为一个李代数，D 为 g 的导子李代数。D 的恒等映射定义了 D 与 g 的半直积 h，称为 g 的全纯代数。将 g 视为 h 的一个理想，将 D 视为 h 的一个子代数。
\end{enumerate}

\begin{enumerate}
\def\labelenumi{(\alph{enumi})}
\item
  证明 \(\mathfrak{g}^*\) 在 \(\mathfrak{g}\) 中的中心化子 \(\mathfrak{h}\) 是集合 \(\operatorname{ad}_{\mathfrak{g}}x - x(x\in \mathfrak{g})\)。
\item
  对于 \(x + d\) 的每个元素 \(\mathfrak{h}\)（\(x \in \mathfrak{g}, d \in \mathfrak{D}\)），记
\end{enumerate}

\[
\theta (x + d) = \mathrm{ad} _ {\mathfrak {g}} x - x + d.
\]

证明 \(\theta\) 是 \(\mathfrak{h}\) 的二阶自同构，并将 \(\mathfrak{g}\) 映到 \(\mathfrak{g}^*\)，因而可以将 \(\mathfrak{h}\) 看作 \(\mathfrak{g}^*\) 的全纯代数。

\begin{enumerate}
\def\labelenumi{(\alph{enumi})}
\setcounter{enumi}{2}
\item
  对于每个理想 \(\mathfrak{k} \neq \{0\}\)，若其为 \(\mathfrak{h}\) 的理想，则 \(\mathfrak{k} \cap (\mathfrak{g} + \mathfrak{g}^*) \neq \{0\}\)。（若 \(\mathfrak{k} \cap \mathfrak{g} = \{0\}\)，则 \(\mathfrak{k} \subset \mathfrak{g}^*\)。）
\item
  若 g 是交换的，则 h 的每个非零理想 \(t \neq \{0\}\) 都包含 g。若 K 是一个域且 \(\dim_{K} g < +\infty\)，推出当 \(t \neq g\) 时，h 不可能是 t 的全纯代数（使用 (b)）。
\end{enumerate}

\begin{enumerate}
\def\labelenumi{\arabic{enumi}.}
\setcounter{enumi}{16}
\tightlist
\item
  设 K 是一个域。令 T 为一个不定元，且 \(L = K((T))\)；L 具有一个规范赋值（形式幂级数的阶），从而成为完备赋值域。
\end{enumerate}

\begin{enumerate}
\def\labelenumi{(\alph{enumi})}
\item
  设 \(\mathrm{V}\) 是 \(\mathbf{K}\) 上的有限维向量空间。则 \(\mathrm{V}_{(\mathbb{L})}\) 规范地具有完备可度量化的拓扑向量空间结构（《拓扑向量空间》，第 I 章第 2 节第 3 号，定理 2）。设 \((x_{p})_{p\in \mathbb{Z}}\) 是 \(\mathrm{V}\) 中元素组成的族，且 \(x_{p} = 0\) 对于小于某个整数的 \(p\) 成立。则级数 \(\sum_{-\infty}^{+\infty}x_p\mathrm{T}^p\) 在 \(\mathrm{V}_{(\mathrm{L})}\) 中收敛，并且 \(\mathrm{V}_{(\mathrm{L})}\) 的每个元素都能唯一地以这种方式表示。（取 V 的一个基。）
\item
  向量空间 \(\mathcal{L}_{\mathrm{L}}(\mathrm{V}_{(\mathrm{L})}) = \mathcal{L}_{\mathrm{K}}(\mathrm{V})_{(\mathrm{L})}\) 具有规范的完备可度量化拓扑向量空间结构。 \(\mathrm{V}_{(\mathrm{L})}\) 的每个自同态都能唯一地表示为 \(\sum_{-\infty}^{+\infty}u_p\mathrm{T}^p\) 的形式，其中 \(u_{p}\in \mathcal{L}_{\mathrm{K}}(\mathrm{V})\)，且当 \(u_{p}\) 小于某个整数时 \(p\) 为零。
\item
  如果 \(\mathbf{K}\) 的特征为 0 且 \(u\in \mathcal{L}_{\mathrm{K}}(\mathrm{V})\)，则记作
\end{enumerate}

\[
e ^ {\mathrm{Tu}} = 1 + \frac {1}{1 !} u \mathrm{T} + \frac {1}{2 !} u ^ {2} \mathrm{T} ^ {2} + \dots \in \mathcal {L} _ {\mathrm{L}} (\mathrm{V} _ {(\mathrm{L})}).
\]

若 \(x \in V\) 满足 \(ux = 0\)，则 \(e^{Tu}.x = x\)。

\begin{enumerate}
\def\labelenumi{(\alph{enumi})}
\setcounter{enumi}{3}
\item
  若 W 是 K 上另一个有限维向量空间，且 \(v \in \mathcal{L}_{\mathrm{K}}(\mathrm{W})\)，则 \(e^{\mathrm{T}(u \otimes 1 + 1 \otimes v)} = e^{\mathrm{T}u} \otimes e^{\mathrm{T}v}\)。
\item
  由 (c) 和 (d) 推出，若 V 具有一个（不必结合的）代数结构且 u 是 V 的导子，则 \(e^{Tu}\) 是代数 \(V_{(L)}\) 的自同构。
\end{enumerate}

¶ 18. 设 H 是一个复 Hilbert 空间，Y 和 Z 是 H 上的连续厄米算子，且 B = {[}Z, Y{]}。设 Y 与 B 可交换。

\begin{enumerate}
\def\labelenumi{(\alph{enumi})}
\item
  证明 \([Z, Y^n] = nBY^{n-1}\)（对 \(n\) 作归纳）。
\item
  设 \(f\) 是实变量的连续可微函数。由 (a) 推出 \([Z, f(Y)] = Bf'(Y)\)。
\item
  由 (b) 推出 \(B = 0\)。（证明若 \(B \neq 0\)，可以选取 \(f\) 使得 \(\| f(Y) \| \leqslant 1\)，而 \(\| Bf'(Y) \|\) 任意大。）
\item
  设 g 是 H 上满足 \(T^{*} = -T\) 的连续算子 T 所成的李代数。由 (c) 推出，条件 \(Y \in g\)、\(Z \in g\)、\([[Z, Y], Y] = 0\) 蕴含 \([Z, Y] = 0\)。
\end{enumerate}

\begin{enumerate}
\def\labelenumi{\arabic{enumi}.}
\setcounter{enumi}{18}
\tightlist
\item
  \begin{enumerate}
  \def\labelenumii{(\alph{enumii})}
  \tightlist
  \item
    设 \(\mathbf{L}\) 是 \(\mathbf{K}\) 上的结合代数。给定 \(k\) 个元素 \(x_{1},\ldots ,x_{k}\) 属于 \(\mathbf{L}\)，记：
  \end{enumerate}
\end{enumerate}

\[
f _ {k} (x _ {1}, \dots , x _ {k}) = \sum_ {\sigma \in \mathfrak {S}_{k}} x _ {\sigma (1)} \dots x _ {\sigma (k)}.
\]

在代数 \(L \otimes_{K} K[T_{1}, \ldots, T_{k}]\) 中，其中 \(T_{i}\) 是不定元，\(f_{k}(x_{1}, \ldots, x_{k})\) 是 \(T_{1} \ldots T_{k}\) 中 \((x_{1}T_{1} + \cdots + x_{k}T_{k})^{k}\) 的系数；它也是下式中 \(T_{1} \ldots T_{k}\) 的系数：

\[
\sum_ {j = 0} ^ {k - 1} \left(x _ {1} \mathrm{T} _ {1} + \dots + x _ {k - 1} \mathrm{T} _ {k - 1}\right) ^ {j} x _ {k} \mathrm{T} _ {k} \left(x _ {1} \mathrm{T} _ {1} + \dots + x _ {k - 1} \mathrm{T} _ {k - 1}\right) ^ {k - j - 1}.
\]

  给定 \(x, y\)（它们是 \(\mathbf{L}\) 中的两个元素），定义 \(g_{i,k-i}(x,y)\) 为 \(\mathrm{T}_1^i\mathrm{T}_2^{k - i}\) 中 \((x\mathrm{T}_1 + y\mathrm{T}_2)^k\) 的系数。证明

\[
i! (k - i)! g _ {i, k - i} (x, y) = f _ {k} \left(t _ {1}, \dots , t _ {k}\right)
\]

其中当 \(t_h = x\) 时 \(h \leqslant i\)，当 \(t_h = y\) 时 \(h \geqslant i + 1\)。

\begin{enumerate}
\def\labelenumi{(\alph{enumi})}
\setcounter{enumi}{1}
\tightlist
\item
  若将 L 视为 K 上的李代数，则在 \(\mathcal{L}_{\mathrm{K}}(\mathrm{L})\) 中：
\end{enumerate}

\[
(\operatorname{ad} x) ^ {n} = \sum_ {i = 0} ^ {n} (- 1) ^ {n - i} \binom {n} {i} \mathrm{L} _ {x} ^ {i} R _ {x} ^ {n - i},
\]

其中 \(\mathbf{L}_x\)（分别为 \(R_{x}\)）是乘法 \(y\mapsto xy\)（分别为 \(y\mapsto yx\)）。

推出，若 K 的特征为 p（p 为素数），则：

\begin{enumerate}
\def\labelenumi{(\arabic{enumi})}
\tightlist
\item
\end{enumerate}

\[
(\operatorname{ad} x) ^ {p}. y = (\operatorname{ad} x ^ {p}). y\tag{2}
\]

\[
(\operatorname{ad} x) ^ {p - 1}. y = \sum_ {i = 0} ^ {p - 1} x ^ {i} y x ^ {p - 1 - i}.
\]

由 (2) 推出：

\[
f _ {p - 1} (\operatorname{ad} x _ {1}, \dots , \operatorname{ad} x _ {p - 1}). y = f _ {p} (x _ {1}, \dots , x _ {p - 1}, y)
\]

以及：

\[
g _ {i, p - 1 - i} (\mathrm{ad} x, \mathrm{ad} y). y = (p - i) g _ {i, p - i} (x, y).
\]

从而对 L 中任意两个元素 x、y：

\[
(x + y) ^ {p} = x ^ {p} + y ^ {p} + \Lambda_ {p} (x, y)\tag{3}
\]

其中：

\[
\Lambda_ {p} (x, y) = \sum_ {i = 1} ^ {p - 1} (p - i) ^ {- 1} g _ {i, p - 1 - i} (\mathrm{ad} x, \mathrm{ad} y). y
\]

属于由 x 和 y 生成的 L 的李子代数（“Jacobson 公式”）。

\begin{enumerate}
\def\labelenumi{\arabic{enumi}.}
\setcounter{enumi}{19}
\tightlist
\item
  设 g 是环 K 上满足 \(pK = \{0\}\) 的李代数（p 为素数）。g 到自身的映射 \(x \mapsto x^{[p]}\) 称为 p-映射，如果它满足关系式
\end{enumerate}

\[
\begin{array}{r l} & {\operatorname{ad} x ^ {[ p ]} = (\operatorname{ad} x) ^ {p}} \\ & {(\lambda x) ^ {[ p ]} = \lambda^ {p} x ^ {[ p ]}} \\ & {(x + y) ^ {[ p ]} = x ^ {[ p ]} + y ^ {[ p ]} + \Lambda_ {p} (x, y) \quad (\mathrm{cf.Exercise19(b)})} \end{array}
\]

对于 \(x, y\)（属于 \(\mathfrak{g}\)）以及 \(\lambda \in \mathbf{K}\)。李 \(p\)-代数定义在 \(\mathbf{K}\) 上，是赋予其一个 \(\mathbf{K}\) 上的李代数结构和一个 \(p\)-映射所得到的结构。每个 \(\mathbf{K}\) 上的结合代数都是李 \(p\)-代数，其中 \([x, y] = xy - yx\) 且 \(x^{[p]} = x^p\)（练习 19）。映射 \(u\) 将李 \(p\)-代数 \(\mathfrak{g}\) 映入李 \(p\)-代数 \(\mathfrak{g}'\)；称为 \(p\)-同态，只要 \(u\) 是李代数同态且 \(u(x^{[p]}) = (u(x))^{[p]}\)。证明，在李 \(p\)-代数 \(\mathfrak{g}\) 中，每个 \(p\)-映射都具有 \(x \mapsto x^{[p]} + f(x)\) 的形式，其中 \(f\) 是将 \(\mathfrak{g}\) 映入其中心的映射，并且关于自同态 \(\lambda \mapsto \lambda^p\)（作用于 \(\mathbf{K}\)）是半线性的。更一般地，若 \(u\) 是从 \(\mathfrak{g}\) 到李 \(p\)-代数的同态，其中 \(\mathfrak{g}', u(x^{[p]}) - (u(x))^{[p]}\) 属于 \(\mathfrak{g}'\) 的中心。

\begin{enumerate}
\def\labelenumi{\arabic{enumi}.}
\setcounter{enumi}{20}
\tightlist
\item
  \begin{enumerate}
  \def\labelenumii{(\alph{enumii})}
  \tightlist
  \item
    设 \(\mathbf{K}\) 是特征为 \(p > 0\) 的域，\(\mathbf{L}\) 是 \(\mathbf{K}\) 上不必结合的代数。证明 \(\mathbf{L}\) 的导子构成一个李 \(p\)-代数，其 \(p\)-映射为 \(\mathbf{D} \mapsto \mathbf{D}^p\)。
  \end{enumerate}
\end{enumerate}

\begin{enumerate}
\def\labelenumi{(\alph{enumi})}
\setcounter{enumi}{1}
\item
  设 \(\mathbf{L}\) 为代数 \(\mathrm{K}[\mathrm{X}] / (\mathrm{X}^p)\)，\(\mathfrak{w}\) 为 \(p\) 的导子所成的 \(\mathbf{L}\)-代数。令 \(\mathrm{D}_i\)（\(0 \leqslant i \leqslant p - 1\)）表示满足 \(\mathbf{L}\) 的 \(\mathrm{D}_i(\mathbf{X}) = \mathbf{X}^i\) 的导子。证明 \(\mathrm{D}_i\) 构成 \(\mathfrak{w}\) 在 \(\mathbf{K}\) 上的一组基，并且当 \([\mathrm{D}_i, \mathrm{D}_j] = (j - i)\mathrm{D}_{i + j - 1}\) 时 \(i + j \leqslant p\)，当 \([\mathrm{D}_i, \mathrm{D}_j] = 0\) 时 \(i + j > p\)，且 \(\mathrm{D}_i^p = 0\)；但 \(i = 1\) 时例外，此时 \(\mathrm{D}_1^p = \mathrm{D}_1\)。
\item
  证明，若 \(p \geqslant 3\)，李代数 \(\mathfrak{w}\) 只有理想 \(\{0\}\) 和 \(\mathfrak{w}\)。（利用 \(D_i\) 的乘法表，证明 \(\mathfrak{w}\) 的每个非零理想都包含 \(D_{p-1}\) 的一个非零倍数。）
\item
令 \(V_{k}\) 为 w 的子空间，其基由 \(D_{i}\)（指标为 \(\geqslant k\)）构成。证明，若 \(p \geqslant 5\)，则 \(V_k\) 恰为 \(Z \in w\) 中中心化子维数为 \(\geqslant k + 1\) 的元素所成的集合。推出，若 \(p \geqslant 5\)，\(w\) 的自同构群是可解的。
\end{enumerate}

\begin{enumerate}
\def\labelenumi{\arabic{enumi}.}
\setcounter{enumi}{21}
\tightlist
\item
  设 g 是特征为 p 的环 K 上的李 p-群（p 为素数）。令 \(x \mapsto x^{[p]}\) 表示 p-映射。对于 g 的每个子集 E，令 \(E^{p}\) 表示由 \(x^{p}\)（其中 \(x \in E\)）生成的 g 的子模。理想 \(a \subset g\) 若满足 \(a^{p} \subset a\)，则称为 p-理想。
\end{enumerate}

\begin{enumerate}
\def\labelenumi{(\alph{enumi})}
\item
对于理想 \(\mathfrak{a} \subset \mathfrak{g}\)，它为 \(p\)-理想的充要条件是它是某个 \(p\)-同态的核（练习 20）。
\item
两个 \(p\)-理想之和是 \(p\)-理想。\(p\)-理想与 \(p\)-子代数之和是 \(p\)-子代数。
\item
设 \(\mathfrak{h}\) 是 \(\mathfrak{g}\) 的李子代数。最小李 \(p\)-子代数（包含 \(\mathfrak{h}\)）是 \(\mathfrak{h} + \mathfrak{h}^p + \mathfrak{h}^{p^2} + \cdots = \overline{\mathfrak{h}}\)；若记 \(\mathfrak{h}_i = \mathfrak{h} + \mathfrak{h}^p + \cdots + \mathfrak{h}^{p^i}\)，则 \(\mathfrak{h}_i\) 是 \(\overline{\mathfrak{h}}\) 的理想，且 \(\overline{\mathfrak{h}}/\mathfrak{h}\) 是交换的。若 \(\mathfrak{h}\) 是 \(\mathfrak{g}\) 的理想，则 \(\mathfrak{h}_i\) 和 \(\overline{\mathfrak{h}}\) 也是理想，并且后者是最小 \(p\)-理想，包含 \(\mathfrak{h}\)。
\item
  g 的任意子模的正规化子和中心化子都是 g 的 p-子代数。
\end{enumerate}

\begin{enumerate}
\def\labelenumi{\arabic{enumi}.}
\setcounter{enumi}{22}
\tightlist
\item
设 \(\mathfrak{g}\) 是一个有限维交换李 \(p\)-代数，定义在完美域 \(\mathbf{K}\) 上，且其特征为 \(p > 0\)。
\end{enumerate}

\begin{enumerate}
\def\labelenumi{(\alph{enumi})}
\item
  证明 g 唯一地是两个 p-子代数 h、t 的直和，其中 p-映射在 h 上是双射，在 t 上是幂零的（考察 g 上 p-映射的迭代；参见《代数学》，第 VIII 章第 2 节第 2 号引理和第 VII 章第 5 节练习 20 (a)）。h 称为 g 的 p-核。
\item
  设 K 是代数闭域。证明存在一组基 \((e_{i})\)，使 \(\mathfrak{h}\) 中的这些基元满足 \(e_{i}^{p}=e_{i}\)（考察由单个元素生成的 h 的 p-子代数，证明它包含某个满足 \(x^{p}=x\neq0\) 的 x；然后对 \(\mathfrak{h}\) 的维数作归纳）。推出，h 的 p-子代数只有有限多个，并且它们与由 \(F_{p}\) 生成的素域 \(e_{i}\) 上的向量空间的向量子空间一一对应。
\item
证明存在一组基 \((f_{ij})\)，属于 \(\mathfrak{t}\)（\(1 \leqslant i \leqslant k\)、\(1 \leqslant j \leqslant s_i\)，对所有 \(i\)），其中 \(s_i\) 的序列递减，并满足 \(f_{1j}^p = 0\)（当 \(1 \leqslant j \leqslant s_1\) 时），\(f_{ij}^p = f_{i-1,j}\)（当 \(2 \leqslant i \leqslant k\)、\(1 \leqslant j \leqslant s_i\) 时）（方法见《代数学》，第 VII 章第 5 节练习 20 (b)）。
\end{enumerate}

¶ 24. 设 K 为一个（交换）域，具有（任意）特征 p，并设 \(X_{i}, Y_{i}, Z_{i}\) 是 3n 个彼此不同的不定元；为简便起见，分别将系统 \((\mathrm{X}_{i})_{1 \leq i \leq n}, (\mathrm{Y}_{i})_{1 \leq i \leq n}, (\mathrm{Z}_{i})_{1 \leq i \leq n}\) 记为 x、y、z。令 \(A_{x,y,z}\) 为系数在 K 中、关于 3n 个不定元 \(X_{i}, Y_{i}, Z_{i}\) 的形式幂级数代数；将 \(A_{x,y}\)（分别为 \(A_{x}\)）记作 \(A_{x,y,z}\) 的不含 z（分别不含 y、z）的幂级数所成的子代数。\(A_{x,y,z}\)（分别为 \(A_{x,y}, A_{x}\)）中的元素记为 \(u(\mathbf{x}, \mathbf{y}, \mathbf{z})\)（分别记为 \(u(\mathbf{x}, \mathbf{y}), u(\mathbf{x})\)）；

  一组 \((u_{i}(\mathbf{x},\mathbf{y},\mathbf{z}))_{1\leqslant i\leqslant n}\) 由 \(A_{x,y,z}\) 中的 n 个元素组成，记为 \(u(\mathbf{x},\mathbf{y},\mathbf{z})\)；对于只含三个系统 x、y、z 中一个或两个的形式幂级数，也采用类似记号。如果 \(u(\mathbf{x},\mathbf{y},\mathbf{z})\in A_{x,y,z}\)，且 \(\mathbf{f}=(f_{i}),\mathbf{g}=(g_{i}),\mathbf{h}=(h_{i})\) 是由 \(A_{x,y,z}\) 中没有常数项的 n 个形式幂级数组成的三个系统，则用 \(u(\mathbf{f}(\mathbf{x},\mathbf{y},\mathbf{z}),\mathbf{g}(\mathbf{x},\mathbf{y},\mathbf{z}),\mathbf{h}(\mathbf{x},\mathbf{y},\mathbf{z}))\) 表示将 \(f_{i}\) 代入 \(X_{i}, g_{i}\)、将 \(Y_{i}, h_{i}\) 代入 \(Z_{i}\)、并按 \(1\leqslant i\leqslant n\) 的范围在 u 中得到的形式幂级数。对于每个系统 \(\alpha=(\alpha_{1},\ldots,\alpha_{n})\)（它由 n 个满足 \(\geqslant0\) 的整数构成），令 \(x^{\alpha}\) 表示单项式 \(X_{1}^{\alpha_{1}}\ldots X_{n}^{\alpha_{n}}\)；\(y^{\alpha}\) 和 \(z^{\alpha}\) 的定义类似。令 \(\varepsilon_{i}\) 表示这样的系统 \(\alpha\)，其满足 \(\alpha_{j}=0\)，条件为 \(j\neq i, \alpha_{i}=1\)。我们将系统 \((0,\ldots,0)\) 记为 e；它由 n 个 \(A_{x,y,z}\) 中的元素组成。

\begin{enumerate}
\def\labelenumi{(\alph{enumi})}
\item
  K 上的形式群律（或者滥用术语称 K 上的形式群）是维数为 n 的 \(G = \mathbf{f}(\mathbf{x}, \mathbf{y})\) 中 n 个元素的系统 \(A_{x,y}\)，满足以下性质：(1) \(\mathbf{f}(\mathbf{x}, \mathbf{f}(\mathbf{y}, \mathbf{z})) = \mathbf{f}(\mathbf{f}(\mathbf{x}, \mathbf{y}), \mathbf{z})\)；(2) \(\mathbf{f}(\mathbf{e}, \mathbf{y}) = \mathbf{y}\)、\(\mathbf{f}(\mathbf{x}, \mathbf{e}) = \mathbf{x}\)。证明必有 \(f_{i}(\mathbf{x}, \mathbf{y}) = \mathrm{X}_{i} + \mathrm{Y}_{i} + g_{i}(\mathbf{x}, \mathbf{y})\)，其中 \(g_{i}\) 只含次数 \(\geqslant 2\) 的单项式，且每个单项式都至少含有一个 \(X_{j}\) 和一个 \(Y_{j}\)。证明存在唯一一个由 \(\mathbf{h}(\mathbf{x})\) 中 n 个元素组成的系统 \(A_{x}\)，使得 \(\mathbf{f}(\mathbf{x}, \mathbf{h}(\mathbf{x})) = \mathbf{f}(\mathbf{h}(\mathbf{x}), \mathbf{x}) = \mathbf{e}\)（《代数学》，第 IV 章第 5 节第 9 号命题 10）。若 \(\mathbf{f}(\mathbf{y}, \mathbf{x}) = \mathbf{f}(\mathbf{x}, \mathbf{y})\)，则称 G 为交换的。
\item
对于所有 \(u \in \mathbf{A}_{\mathbf{x}}\)，令 \(L_{\mathbf{y}}u\) 表示 \(u(\mathbf{f}(\mathbf{y},\mathbf{x}))\) 中的元素 \(\mathbf{A}_{\mathbf{x},\mathbf{y}}\)。\(\mathbf{A}_{\mathbf{x}}\) 的导子 D 可规范地延拓为 \(\mathbf{A}_{\mathbf{x},\mathbf{y}}\) 的导子（同样记为 D，这里滥用记号），约定 \(\mathrm{D}(\mathrm{Y}_i) = 0\) 对所有 \(i\) 成立（《代数学》，第 IV 章第 5 节第 8 号命题 6）。若 \(L_{\mathbf{y}}D = DL_{\mathbf{y}}\)，则称 D 关于所讨论的形式群 G 左不变。记 \(\mathrm{D}^{(\mathrm{e})}\) 为从 \(\mathbf{A}_{\mathbf{x},\mathbf{y}}\) 到 \(\mathbf{A}_{\mathbf{y}}\) 的线性映射，其定义为 \(\mathrm{D}^{(\mathrm{e})}(u) = (\mathrm{D}u)(\mathbf{e},\mathbf{y})\)；它将 \(\mathbf{A}_{\mathbf{x}}\) 映入 K，并由其在 \(\mathbf{A}_{\mathbf{x}}\) 上的限制确定。证明 D 左不变的充要条件是对所有 \(v \in \mathbf{A}_{\mathbf{x}}\)，都有 \(\mathrm{D}^{(\mathrm{e})}(L_{\mathbf{y}}v) = (\mathrm{D}v)(\mathbf{y})\)。
\end{enumerate}

令 \(D_{i}\) 为 \(\partial/\partial X_{i}\) 上的导子 \(A_{x}\)（\(1 \leqslant i \leqslant n\)）；存在唯一一个左不变导子 \(T_{i}\) 满足 \(T_{i}^{(e)} = D_{i}^{(e)}\)。证明 \(T_{i}\) 在 K 上线性无关，并且每个左不变导子都是系数属于 K 的 \(T_{i}\) 的线性组合。对于括号 {[}D, D'{]} 和 p-映射 \(D \mapsto D^{p}\)（当 p \textgreater{} 0 时），推出左不变导子所成的集合 g 是一个李代数（当 p \textgreater{} 0 时是李 p-代数），称为形式李群 G 的李代数。

\begin{enumerate}
\def\labelenumi{(\alph{enumi})}
\setcounter{enumi}{2}
\tightlist
\item
  证明，对每个形式幂级数 \(u \in \mathbf{A}_{\mathbf{x}}\) :
\end{enumerate}

\[
u (\mathbf {f} (\mathbf {x}, \mathbf {y})) = u (0) + \sum_ {i = 1} ^ {n} Y _ {i} T _ {i} u + o _ {2} (\mathbf {x}, \mathbf {y})\tag{1}
\]

其中级数 \(o_{2}\)中的所有单项式次数均为 \(\geqslant2\)，关于 \(Y_{j}\)。推出

\[
\begin{array}{r l} {u (\mathbf {f} (\mathbf {x}, \mathbf {f} (\mathbf {y}, \mathbf {z}))) = u (0) + \sum_ {i = 1} ^ {n} (\mathrm{Y} _ {i} + \mathrm{Z} _ {i}) \mathrm{T} _ {i} u} & \\ {+ \sum_ {i, j} \mathrm{Y} _ {i} \mathrm{Z} _ {j} ((\mathrm{T} _ {i} \mathrm{T} _ {j}) (u)) + o _ {3} (\mathbf {x}, \mathbf {y}, \mathbf {z})} & \end{array}\tag{2}
\]

其中级数 \(o_{3}\)中的所有单项式次数均为 \(\geqslant3\)，关于集合 \(Y_{i}\) 和 \(Z_{i}\)。推出：

\[
u (\mathbf {f} (\mathbf {x}, \mathbf {y})) - u (\mathbf {f} (\mathbf {y}, \mathbf {x})) = \sum_ {i <   j} (X _ {i} Y _ {j} - X _ {j} Y _ {i}) ([ T _ {i}, T _ {j} ] ^ {(e)} (u)) + o _ {3} ^ {\prime} (\mathbf {x}, \mathbf {y})\tag{3}
\]

其中 \(o_3'\)中的所有项的总次数均为 \(\geqslant 3\)。证明，若 \(G\) 是交换的，\(g\) 是交换的。

\begin{enumerate}
\def\labelenumi{(\alph{enumi})}
\setcounter{enumi}{3}
\tightlist
\item
  设 G' 是另一个维数为 m 的形式群，其群律为 f'。 从 G 到 G' 的形式同态是一个系统 \(\mathbf{F} = (\mathrm{F}_{j}(\mathbf{x}))_{1 \leqslant j \leqslant m}\)，其元素属于 \(\mathbf{A}_{\mathbf{x}}\)，满足 \(\mathbf{f}'(\mathbf{F}(\mathbf{x}), \mathbf{F}(\mathbf{y})) = \mathbf{F}(\mathbf{f}(\mathbf{x}, \mathbf{y}))\) 。对于每个左不变导子 \(D \in g\) ，证明 \(v \mapsto D(v \circ F)\) 是属于李代数 \(g'\) 的左不变导子。 若记为 \(\mathbf{F}^*(D)\) , \(\mathbf{F}^*\)  是同态（当 p \textgreater{} 0 时为 p-同态），从 g 到 \(g'\) 。若 \((\mathrm{T}_{1})_{1 \leqslant i \leqslant n}, (\mathrm{T}'_{j})_{1 \leqslant j \leqslant m}\) 是 g 和 \(g'\) 的基，满足 \(T_{i}^{(e)} = D_{i}^{(e)}\)，且 \(T'_{j}{}^{(e)} = D'_{j}{}^{(e)}\) ，证明 \(F^*\) 关于这些基的矩阵为矩阵 \((\partial F_{j}/\partial X_{i})_{0}\)（\(\partial F_{j}/\partial X_{i}\)  的常数项；其中 i 是行的指标，j 是列的指标）。
\end{enumerate}

\begin{enumerate}
\def\labelenumi{\arabic{enumi}.}
\setcounter{enumi}{24}
\tightlist
\item
  在 \(n\) 个变量上，域 \(\mathbf{K}\)，记为 \(\mathbf{GL}(n, \mathbf{K})\)（无混淆时），是维数为 \(n^2\)，其群律由 \(f_{ij}(\mathbf{u}, \mathbf{v}) = -\delta_{ij} + \sum_{k=1}^{n} (\delta_{ik} + \mathrm{U}_{ik})(\delta_{kj} + \mathrm{V}_{kj})\)（\(\delta_{ij}\) 为克罗内克指标；一般定义中练习 24 的 \(3n^2\) 个不定元在这里记为 \(\mathrm{U}_{ij}, \mathrm{V}_{ij}, \mathrm{W}_{ij}\)，两个指标从 1 变化到 \(n\)）。证明，若 \(\mathrm{D}_{ij} = \partial/\partial \mathrm{U}_{ij}\)，满足该条件的左不变导子 \(\mathrm{X}_{ij}\) 满足 \(\mathrm{X}_{ij}^{(e)} = \mathrm{D}_{ij}^{(e)}\) 由下式给出
\end{enumerate}

\[
\mathrm{X} _ {i j} = (1 + \mathrm{U} _ {i i}) \mathrm{D} _ {i j} + \sum_ {k \neq i} \mathrm{U} _ {k i} \mathrm{D} _ {k j}.
\]

 的李代数 \(\mathbf{GL}(n,\mathbf{K})\) 被识别为 \(\mathfrak{gl}(n,\mathbf{K})\)，将 \(X_{ij}\) 与规范基中的元素 \(E_{ij}\)（使用练习 24 的公式 (3)）。若 K 的特征为 p \textgreater{}，\(X_{ii}^{p} = X_{ii}\)，且 \(X_{ij}^{p} = 0\) 当 \(i \neq j\)。

¶ 26. (a) 设 K 为特征为 \(\neq 2\)  的域，E 为 \(n\)  维 K 向量空间，\(\Phi\) 是 E 上的非退化对称双线性型。设 E 关于 \(\Phi\) 具有一组正交基，令 \(R\) 表示 \(\Phi\) 关于该基的每个矩阵 \(U\) 属于正交群 \(\mathbf{G}(\Phi)\) 且满足 \(\det(I + U) \neq 0\) 可写成

\[
U = (I - R ^ {- 1} S) ^ {- 1} (I + R ^ {- 1} S),
\]

  其中 \(S = R(U - I)(U + I)^{-1}\) 是交错矩阵，满足 \(\det(I - R^{-1}S) \neq 0\)；反之，对每个矩阵 \(S\) 满足这些条件的矩阵 \(U = (I - R^{-1}S)^{-1}(I + R^{-1}S)\) 是 \(\mathbf{G}(\Phi)\) 的矩阵，满足 \(\det(I + U) \neq 0\) 。令 \(S_{ij}, S_{ij}', S_{ij}''\) 为 \(3n(n - 1)/2\) 个不定元（ \(1 \leqslant i < j \leqslant n\) )

并令 L 为包含形式幂级数环 \(K[[S_{ij}, S'_{ij}, S''_{ij}]]\)。若以 \(S, S', S''\) 表示矩阵

\[
\sum_ {i <   j} \mathrm{S} _ {i j} (E _ {i j} - E _ {j i}), \qquad \sum_ {i <   j} \mathrm{S} _ {i j} ^ {\prime} (E _ {i j} - E _ {j i}), \qquad \sum_ {i <   j} \mathrm{S} _ {i j} ^ {\prime \prime} (E _ {i j} - E _ {j i}),
\]

则 \(\operatorname{det}(I - R^{-1}S) \neq 0\)，以及对 \(S'\) 和 \(S''\) 的类似结论也成立；若记

\[
U = (I - R ^ {- 1} S) ^ {- 1} (I + R ^ {- 1} S), \quad U ^ {\prime} = (I - R ^ {- 1} S ^ {\prime}) ^ {- 1} (I + R ^ {- 1} S ^ {\prime}),
\]

则还有 \(\operatorname{det}(I + UU') \neq 0\) . 记

\[
F (S, S ^ {\prime}) = R \left(U U ^ {\prime} - I\right) \left(U U ^ {\prime} + I\right) ^ {- 1} = \left(f _ {i j} \left(S, S ^ {\prime}\right)\right);
\]

其中 \(f_{ij}(S, S')\) 属于形式幂级数环 \(K[[S_{ij}, S_{ij}']]\)，并将 \(F(S, S')\) 看作系数属于 K 上 n 阶矩阵代数的形式幂级数，可写为：

\[
F (S, S ^ {\prime}) = S + S ^ {\prime} + o _ {2} (S, S ^ {\prime})
\]

其中，在矩阵\(o_2(S, S')\)的元素中，各项的次数关于\(\geqslant 1\)均为\(S_{ij}\)，且次数关于\(\geqslant 1\)均为\(S_{ij}'\)；同理：

\[
U = I + 2 R ^ {- 1} S + o _ {2} ^ {\prime} (S)
\]

 的元素是 \(o_2'(S)\) 为阶数 \(\geqslant 2\)。证明 \(F(S, S')\) 是 \(K\) 上的形式群律，维数为 \(n(n - 1)/2\) ；该形式群称为形式正交群（对应于 \(\Phi\) ），并记为 \(G(\Phi)\)（这里滥用记号）。

\begin{enumerate}
\def\labelenumi{(\alph{enumi})}
\setcounter{enumi}{1}
\tightlist
\item
  若记\(M(S) = (m_{ij}(S)) = (I - R^{-1}S)^{-1}(I + R^{-1}S)\) ，\(M\) 是从 \(\mathbf{G}(\Phi)\) 到 \(\mathbf{GL}(n, \mathbf{K})\)  的形式同态。令 \(\mathfrak{o}(\Phi)\) 为 \(\mathbf{G}(\Phi)\)，并令 \((\mathrm{T}_{ij})_{i < j}\) 表示该代数的一组基，满足 \(\mathrm{T}_{ij}^{(e)} = \mathrm{D}_{ij}^{(e)}\)（参见练习 24 (d)）；若元素 \(D = \sum_{i < j} c_{ij} T_{ij}\) 属于 \(\mathfrak{o}(\Phi)\) 被识别为矩阵 \(D = \sum_{i < j} c_{ij}(E_{ij} - E_{ji})\) ，证明（按练习 24 (d) 的记号，并使用练习 25 中将...的李代数与 \(\mathbf{GL}(n, \mathbf{K})\) 识别为 \(\mathfrak{gl}(n, \mathbf{K})\)）:
\end{enumerate}

\[
M ^ {*} (D) = 2 R ^ {- 1} D.
\]

  推出，\(M^*\) 是 \(\mathfrak{o}(\Phi)\) 到由 \(\mathfrak{gl}(n, \mathbf{K})\) 所成的子代数，矩阵 \(X\) 满足 \({}^{t}XR + RX = 0\)（对每个有序对 \((a, b)\) 个元素 \(E\) ，将其识别为单列矩阵，考察 \(\Phi(a + M(S).a, b + M(S).b) = {}^{t}a.(I + {}^{t}M(S))R(I + M(S)).b\) 作为取常数项后的形式幂级数，并利用由 \(X \in \mathfrak{gl}(n, \mathbf{K})\) 满足 \({}^{t}XR + RX = 0\) 的子代数维数为 \(n(n - 1)/2\)）。

\begin{enumerate}
\def\labelenumi{(\alph{enumi})}
\setcounter{enumi}{2}
\tightlist
\item
  类似地定义任意
\end{enumerate}

（交换）域上，并证明其李代数被识别为由 \(\mathfrak{gl}(2n,\mathbf{K})\) 所成的子代数，其中矩阵 X 满足

\[
{ } ^ { t } X A + A X = 0 , \quad \text {   where   } \quad A = \left( \begin{array} { c c } 0 & I _ { n } \\ - I _ { n } & 0 \end{array} \right) .
\]

\begin{enumerate}
\def\labelenumi{\arabic{enumi}.}
\setcounter{enumi}{26}
\tightlist
\item
  设 K 为特征为 p \textgreater{} 为 0，G 为由 \(f_{1}(\mathbf{x}, \mathbf{y}) = \mathrm{X}_{1} + \mathrm{Y}_{1} + \mathrm{X}_{1}\mathrm{Y}_{1}, f_{2}(\mathbf{x}, \mathbf{y}) = \mathrm{X}_{2} + \mathrm{Y}_{2}(1 + \mathrm{X}_{1}^{p})\)。证明 G 不交换，但其李代数是交换的。
\end{enumerate}

